\documentclass[10pt]{article}
\usepackage[margin=1in]{geometry}
\usepackage{amsmath,amssymb,amsthm}
\usepackage{booktabs}
\usepackage{enumitem}
\usepackage{hyperref}
\setlist{itemsep=1pt,topsep=2pt}

\newtheorem{hypothesis}{Hypothesis}
\newtheorem{finding}{Finding}
\newtheorem{proposition}{Proposition}
\newtheorem{remark}{Remark}
\newcommand{\E}{\mathbb{E}}
\newcommand{\R}{\mathbb{R}}
\newcommand{\ip}[2]{\langle #1, #2 \rangle}

\title{\vspace{-2em}The Bilinear GCIQL Critic Is Not a Learned Distance Function\\
{\large Evidence from OGBench cube manipulation, and a mathematical account of what the
critic learns instead}}
\author{}
\date{August 2026}

\begin{document}
\maketitle
\vspace{-3em}

\paragraph{Setting.} All experiments are offline goal-conditioned RL on OGBench
\texttt{cube-single-play-v0} (state-based single-cube manipulation, one million
transitions of scripted, task-agnostic play data), 3 seeds $\times$ 500{,}000 gradient
steps per method. The policy is a flow-matching behavior-cloning actor plus a critic-value
ascent term; a control actor with the critic term removed isolates the critic's
contribution. Methods compared: goal-conditioned IQL with a standard MLP critic
(\textbf{GCIQL}), the same objective with a bilinear critic (\textbf{bilinear GCIQL}),
contrastive RL (\textbf{CRL}), and quasimetric RL (\textbf{QRL}) as a reference that
enforces distance structure by architecture.

\section{The hypothesis under test}

The bilinear critic is $\hat Q(s,a,g)=\ip{\varphi(s,a)}{\psi(g)}/\sqrt d$ with value
$\hat V(s,g)=\ip{\varphi(s)}{\psi(g)}/\sqrt d$, trained by goal-conditioned expectile TD
(expectile $0.9$, discount $\gamma=0.99$) on goal-relabeled data. Let $d^*(s,g)$ be the
minimal number of steps from $s$ to $g$ supported by the dataset. Any value function
induces an estimated distance $\hat d(s,g) = \log_\gamma\!\big(1+(1-\gamma)\hat V\big)$,
a monotone transform of $\hat V$; every measurement below is rank-based and therefore
identical for $\hat V$ and $\hat d$.

\begin{hypothesis}[The distance hypothesis]\label{hyp:metric}
Bilinear GCIQL outperforms GCIQL ($92.7\pm1.4$ vs.\ $73.2\pm1.3$ success rate; the
behavior-cloning control reaches $8$--$10$) because its $\hat d$ approximates the true
distance $d^*$ more faithfully, and this faithfulness is what improves the policy.
\end{hypothesis}

If the hypothesis holds, the bilinear critic must (a) rank pairs by temporal distance at
least as well as GCIQL on held-out data, (b) agree with the methods that estimate
distances well by construction (QRL, CRL), and (c) improve its distance estimates over
the training window in which its success rate improves.

\section{The evidence: all three implications fail}

\begin{finding}[GCIQL estimates temporal distance better, and loses]\label{f:p1}
Spearman rank correlation between $\hat d(s,g)$ and the time separation of $(s,g)$, over
$10^5$ held-out same-trajectory pairs (second number: partial correlation controlling for
Euclidean distance between raw observations):
\[
\text{GCIQL}:\; 0.63\,/\,0.54 \qquad>\qquad \text{bilinear GCIQL}:\; 0.54\,/\,0.41 .
\]
The critic with the \emph{worse} distance estimate wins by 20 success points. Distance
faithfulness on the data distribution is therefore neither sufficient (GCIQL) nor
necessary (bilinear GCIQL) for policy improvement.
\end{finding}

\begin{finding}[The bilinear critic disagrees with every good distance
estimator]\label{f:rsa}
Compute each method's $2000\times2000$ matrix of pairwise estimated distances on a fixed
state sample, and correlate the matrices between methods (Spearman on the upper
triangles). GCIQL--QRL correlate at $0.88$ and GCIQL--CRL at $0.54$: the three
distance-flavored methods agree with one another. The bilinear critic agrees with none of
them (at most $0.21$) and is instead closest to plain Euclidean distance between raw
observations ($0.45$--$0.52$). If it had learned the temporal distance, it would agree
with the methods that demonstrably do; instead, those methods underperform it by 20--60
success points.
\end{finding}

\begin{finding}[The similarity structure is an artifact of the architecture, not of
learning]\label{f:cos}
After training, the bilinear output is numerically a scaled cosine similarity: the
correlation between $\ip{\varphi}{\psi}$ and $\cos\angle(\varphi,\psi)$ is $0.996$, and
the embedding norms vary by under $9\%$ across inputs. But the \emph{untrained} network
already shows correlation $0.994$: layer normalization concentrates the norms at
initialization. ``It learned a cosine similarity'' is a property of the architecture and
cannot explain an advantage acquired during training.
\end{finding}

\begin{finding}[Distance quality plateaus while performance keeps rising]\label{f:time}
Between 100k and 500k gradient steps, the rank correlation of Finding~\ref{f:p1} moves
only $0.37\to0.41$ while success moves $65\to93$. The quantity the hypothesis says is
causal is flat over the window in which the effect it must explain triples.
\end{finding}

\begin{finding}[Where the two critics really differ: states visited at evaluation
time]\label{f:p2}
Roll out each trained policy on the evaluation tasks (deterministic actions, 3 seeds
$\times$ 50 episodes) and record, for each episode, the Spearman correlation between
$\hat V(s_t, g)$ and the timestep $t$ --- for a sound value function, value should rise
as the policy approaches its goal. Score \emph{both} critics on \emph{both} policies'
trajectories (median per episode, mean$\pm$sd over seeds):
\begin{center}
\begin{tabular}{lcc}
\toprule
 & trajectories of the bilinear-GCIQL & trajectories of the GCIQL\\
 & policy (success $0.95$) & policy (success $0.73$)\\
\midrule
bilinear-GCIQL critic & $0.983\pm0.003$ & $0.871\pm0.060$\\
GCIQL critic          & $\mathbf{0.355\pm0.239}$ & $0.883\pm0.015$\\
\bottomrule
\end{tabular}
\end{center}
On held-out dataset pairs GCIQL ranks better (Finding~\ref{f:p1}); on the states actually
visited by the \emph{successful} policy at evaluation time, its value estimates are close
to uninformative (on one task the median correlation is $-0.80$: value falls as the policy
approaches the goal), while the bilinear critic's value rises almost perfectly
monotonically. Each critic is coherent along its own policy's trajectories (diagonal);
only the bilinear critic remains coherent on the \emph{other} policy's states. What
separates the two critics is not distance quality on the data distribution but coherence
on state--goal pairs \emph{outside} it --- exactly the pairs evaluation interrogates.
\end{finding}

\begin{remark}[Verdict]
Implication (a) fails by Finding~\ref{f:p1}, (b) by
Findings~\ref{f:rsa}--\ref{f:cos}, (c) by Finding~\ref{f:time}, and
Finding~\ref{f:p2} shows the real discriminator lies elsewhere. The bilinear critic also
freely violates the defining axioms of a distance where it wins --- its estimates are
asymmetric and violate the triangle inequality at rates comparable to GCIQL's. Whatever
it learns, calling it a distance function explains none of the observed orderings.
\end{remark}

\section{What the critic is plausibly learning instead}\label{sec:instead}

The surviving account: the bilinear critic learns a \textbf{bounded, low-rank
approximation of the optimal goal-conditioned value function}, and its advantage comes
from how that function class behaves under policy optimization --- not from any distance
structure. The argument has three steps, each tied to a measured quantity.

\paragraph{Setup.} With reward $-1$ per step until the goal is reached, the optimal value
is $V^*(s,g)=\big(\gamma^{d^*(s,g)}-1\big)/(1-\gamma)\in[-100,\,0]$ at $\gamma=0.99$.
Expectile regression with expectile near $1$ pushes $\hat V$ toward this in-sample optimum
on the distribution $\mu$ of relabeled state--goal pairs. Because layer normalization
concentrates $\|\varphi\|$ and $\|\psi\|$, the bilinear class is uniformly bounded:
\[
|\hat V|,\,|\hat Q| \;\le\; \kappa := \|\varphi\|\,\|\psi\|/\sqrt d,
\qquad \kappa\approx 75 \text{ as measured, realized } \hat V\in[-63,\,+6].
\]

\paragraph{Step 1: the bound caps how much the actor can exploit critic error.}
The actor maximizes $\hat Q$ against a behavior-cloning anchor. For any policy $\pi$ the
standard decomposition gives
\[
J(\pi)\;\ge\;J(\pi_{\mathrm{BC}})
\;+\;\E_{s,g}\big[\hat Q(s,\pi)-\hat Q(s,\pi_{\mathrm{BC}})\big]
\;-\;2\sup_{s,a,g}\big|\hat Q-Q^*\big| .
\]
For an unbounded MLP the supremum is uncontrolled precisely outside the data
distribution, where the maximization goes hunting; in the bilinear class it is at most
$\kappa+\sup|Q^*|$. Indeed $\kappa\approx75<100=\sup|V^*|$: the class cannot even
represent the most extreme values, let alone hallucinate beyond them. Errors are forced
to be many, small, and spread out rather than rare and enormous --- matching the
measurement that GCIQL's triangle-inequality violations average $9.5$ steps in magnitude
against the bilinear critic's $4.8$ at comparable rates.

\paragraph{Step 2: low rank forces the critic to fill in unseen state--goal pairs
consistently.}
View the value function on the dataset's states and goals as a matrix
$M_{sg}=\hat V(s,g)$; the bilinear form constrains $\operatorname{rank}(M)\le d$.
Goal relabeling covers each state with \emph{some} goals and each goal with \emph{some}
states, but not all combinations; evaluation asks about combinations never seen together.
A rank-$d$ matrix has only $d(|\mathcal S|+|\mathcal G|-d)$ degrees of freedom and ---
this is the standard matrix-completion fact --- is uniquely determined by a number of
generically-positioned entries proportional to $d(|\mathcal S|+|\mathcal G|)$, up to
logarithmic factors. Once training pins down the two factors on the covered pairs, the
values on the \emph{uncovered} pairs are no longer free: the class extrapolates them
consistently. An MLP critic obeys no such constraint; its values off the training
distribution are limited only by smoothness. This is Finding~\ref{f:p2} in mathematical
form: the bilinear critic stays coherent on another policy's states ($0.87$) where the
MLP critic collapses ($0.36$).

\paragraph{Step 3: what is learned is a low-rank approximation of $V^*$, not a
distance.}
In the idealized limit, expectile TD in the rank-constrained class solves
\[
\min_{\operatorname{rank} M\le d}\ \big\|\,W\odot(M-V^*)\,\big\|_{\mu}^2 ,
\]
a weighted low-rank approximation of the optimal value function --- i.e.\ the critic
retains (a reweighting of) the top-$d$ singular subspace of $V^*$. This predicts, and the
data confirm:
\begin{enumerate}[label=(\roman*)]
\item performance saturates once $d$ exceeds the effective dimension of that subspace:
measured success $0.06/0.62/0.92/0.90/0.93$ at $d=5/10/20/30/50$, saturating between
$d=10$ and $d=20$, while the trained embeddings' effective dimension keeps growing with
$d$ ($4.1/8.0/14.7/19.6/28.0$) --- capacity, not a learned constant;
\item the axioms of a distance (symmetry, triangle inequality) are incidental --- a
truncated low-rank approximation of $V^*$ satisfies neither exactly, and the bilinear
critic indeed violates both while winning;
\item distance measurements on the data distribution (Findings~\ref{f:p1},
\ref{f:rsa}) can prefer GCIQL while the actor prefers the bilinear critic, because such
measurements are blind to both the bound of Step~1 and the off-distribution consistency
of Step~2.
\end{enumerate}

\begin{proposition}[Summary of the account, and how it can fail]
The bilinear critic's advantage $=$ a uniformly bounded ($\kappa<\sup|V^*|$), rank-$d$
approximation of the optimal goal-conditioned value function, extended to unseen
state--goal pairs by the low-rank constraint. Falsifiable consequences, currently under
test: (a) imposing the same bound $\kappa$ on the MLP critic (squashing its output
through $\kappa\tanh(\cdot/\kappa)$) should recover part of the $73\to93$ gap;
(b) removing the layer normalization from the goal embedding, which destroys the uniform
bound at fixed rank, should degrade the bilinear critic; (c) the off-distribution
coherence of Finding~\ref{f:p2} should vanish for $d$ below the saturation dimension.
Any of these failing narrows the account; none of them restores
Hypothesis~\ref{hyp:metric}, which Findings~\ref{f:p1}--\ref{f:p2} refute independently.
\end{proposition}

\end{document}
